\clearpage
\section{AI-assisted research and writing}\label{sec:ai-use}
The authors proposed the intermediation question, built the original rental pipeline and identified the shortage of affordable listings in the pilot. They used OpenAI Codex (GPT-6 family) extensively to develop and challenge proofs, write and check the national replays, analyze administrative records, and draft and revise the text. The record identifies the model family but does not preserve a checkpoint for every earlier interaction. One request was “Can we test whether the prediction-versus-placement tradeoff holds on a real matching dataset with observed preferences, then fold that evidence into the paper?” A later request asked for a reviewer audit of the theory, experiments and connection between the two settings. The resulting revisions distinguished historical enrollment from simulated placements, tested different comparators and application orders, and checked numerical claims using separate implementations. The rental experiment used the saved selections of the separate model identified as \texttt{claude-haiku-5-5}. Its inputs and cache limitations appear in Sections~\ref{sec:twonyc} and~\ref{sec:limits}.

The subsequent design work included AI-assisted construction of the finite-lottery reduction, vacancy-chain algorithm, cutoff comparison and their counterexamples. The published code records the simulated priority inputs, independent small-instance checks, control policies and aggregate outputs.

The institutional revision used AI assistance to reconstruct seat categories and family rules, compare predictions against the initial-offer file, and rerun the menu optimization and assignment checks. The release includes the comparison code and aggregate results.
